\documentclass[11pt]{article}
\usepackage[margin=1in]{geometry}
\usepackage[T1]{fontenc}
\usepackage{xcolor}
\usepackage{enumitem}
\definecolor{revblue}{RGB}{20,60,140}
\newcommand{\point}[1]{\par\medskip\noindent\textbf{\color{revblue}#1}\par\smallskip}
\newcommand{\reply}{\par\noindent\textbf{Response.}\ }
\setlength{\parindent}{0pt}
\setlength{\parskip}{4pt}
\begin{document}

\begin{center}
{\large\textbf{Response to the reviewers}}\\[2pt]
Aggressive feature reduction degrades the selective reliability and clinical net benefit of
clinical risk models\\[2pt]
Scientific Reports submission 2650209c-ff1e-4df4-aeb0-75ef3b35e564 (revised manuscript)
\end{center}

Greetings, and thank you --- both reviews improved this paper materially, and Reviewer~1's first
point changed one of its conclusions. Rather than patch the affected analyses, we re-ran the entire
experimental grid from scratch under a new, fully hash-pinned configuration, extended it with an
exact-$k$ budget arm and the full Diabetes-130 cohort, and re-derived every number in the manuscript
from the regenerated results. The rebuild reproduces the originally submitted aggregates on the
eleven unchanged datasets at machine precision for the tree-based learners
($\le1.2\times10^{-16}$ on any estimate or interval bound) and within $1\times10^{-4}$ for logistic
regression (largest observed deviation $5.1\times10^{-5}$; a documented solver-version sensitivity,
and the environment is now version-locked). The Holm-corrected $p$-values move correspondingly more
on those logistic cells (at most $3.9\times10^{-3}$), and every one of the 5{,}265 materiality
verdicts is nonetheless identical --- so all changes below are attributable to the revision's
additions, not to instability of the pipeline.

\section*{Reviewer 1}

\point{1. The rounding rule retains at least one feature no matter how aggressive the cut\ldots\
feature count and absolute retained-$k$ are mechanically linked through this rule, so the
correlation may substantially reflect the simpler point that ``reducing to one feature is unsafe.''
I ask the authors to re-relate the harm to the absolute number of retained features (or compare
datasets at matched retained-$k$).}

\reply You were right, and we did both --- and the matched-$k$ experiment overturned our original
reading. First, the reanalysis you asked for: harm correlates with the absolute retained count
\emph{more strongly} than with the feature count (Spearman $\rho=-0.78$ vs.\ $-0.71$; both
leave-one-dataset-out robust with bootstrap intervals excluding zero), and excluding the two
floor-collapsed datasets attenuates the $p$-association to non-significance ($\rho=-0.59$,
$p=0.075$, $n=10$). Second,
the decisive experiment: we added an \emph{exact-$k$ budget arm} ($k\in\{1,2,4,8\}$ for every
dataset, evaluated inside the same repetitions and folds as the fractional budgets, in its own Holm
family so the original comparisons are untouched). At matched $k=1$, \emph{every} dataset is
materially harmed --- including the wide ones: SPECTF ($p=44$), whose own 25\,\% budget costs
$+0.002$ AURC, pays $+0.035$ at $k=1$; WDBC $+0.030$; the full readmission cohort $+0.027$; Statlog
and Cleveland pay $+0.19$ and $+0.18$. The claim you challenged (``datasets with more candidate
features are reliably safer to reduce'') is therefore \emph{retired} in the revision: the apparent
robustness of wide datasets is a property of the retained count their fractional budgets never
leave, dataset structure modulates severity at matched $k$ by an order of magnitude, and the harm is
not attributable to instability of the selection itself. On that last point we report a dissociation
rather than lean on a correlation: Haberman is the second most stable dataset in the study and the
third most harmed, while SPECTF is among the least stable and the least harmed. The correlational
test is given in both directions and as measured --- null across datasets ($\rho=+0.14$, $p=0.66$),
and at the finer (dataset, ranker) level null when pooled ($\rho=-0.047$, $n=36$) but with a weak
within-dataset trend ($\rho=-0.29$, $p=0.19$ by block permutation) whose sign is the one a
noise mechanism would predict. We state plainly that this trend is non-significant, that it is
confounded with ranker quality, and that we do not rest the claim on it; the exact-$k$ arm carries
it instead, since at $k=1$ even datasets with near-deterministic selections are materially harmed.
Section~4.7 is rewritten around this result, the
budget rule's floor and rounding behaviour are now stated explicitly in Methods with their
consequences, and a new table shows that a strict safety certificate (no materially-worse cell on
any axis) passes for exactly one dataset at one budget in this study --- the audit, not any proxy,
is the certificate.

\point{2. I would like more discussion of the precedent for applying feature reduction to
lower-dimensional medical datasets, and an explicit statement that the guidance does not extend to
the high-dimensional regime.}

\reply Both added. Related work now documents the precedent --- most directly, the published
survival study on the very heart-failure cohort we reuse argued that two of its twelve features
suffice, on discrimination-style evidence alone; the Cleveland and Pima datasets were likewise
introduced with compact selected-subset models. The Limitations section now carries an explicit,
bolded scope statement: the guidance does not extend to the high-dimensional regime ($p\gg44$;
genomic, transcriptomic, radiomic, or wide-EHR feature spaces), where selection is a statistical
necessity and retained budgets remain far above the single-digit counts at which we observe the
largest harm.

\point{3. I would like at least one worked example showing where reduction would have flipped a real
treatment decision for a patient, at a specific clinically relevant threshold.}

\reply Added as a new results section, with the cell and threshold fixed \emph{a priori} in a
committed analysis configuration before any flip was counted. On the mammographic-mass data at the
conventional 10\,\% biopsy threshold, reducing to the 25\,\% budget --- which on this dataset is a
single feature --- flips the biopsy decision for 203 of 961 patients (21\,\%) in the majority of the
30 paired models, including 4 patients with malignant masses moved \emph{from biopsy to no biopsy};
a table of exemplar patients with their feature values and per-threshold decision curves is
included, along with sensitivity at $t=0.05$ and $0.15$. At population scale, the same analysis on
the full readmission cohort flips the enrollment decision for 14{,}324 of 101{,}763 encounters
(14.1\,\%), including 1{,}132 true readmissions that lose the intervention. Reliability diagrams
show the mechanism (the reduced model's probabilities compress toward the prevalence), and the
per-patient prediction caches behind these counts are committed, so every number regenerates
without re-running any experiment.

\point{4. I would like the full Diabetes-130 cohort analyzed, at least for the headline analysis.}

\reply Done, beyond the request: the \emph{entire} grid --- all rankers, learners, budgets,
repetitions, and the recalibration control --- now runs on all 101{,}763 eligible encounters, and
every number in the manuscript rests on the full cohort. We also kept the subsample as a robustness
arm and report the agreement explicitly (new Section~4.9): point estimates agree in the typical
case (median absolute difference 0.004), but 20 of 45 materiality verdicts disagree --- in 12 the
full cohort resolves an effect the subsample was underpowered to see, in 8 the subsample asserted
materiality the full cohort refutes, and in 18 of the 45 the subsample reversed the \emph{sign} of
the effect, most consequentially on discrimination ($+0.028$ full versus $-0.028$ subsampled). Your concern was substantively correct, and the manuscript now documents what the
subsample would have missed.

\point{5. As reproducibility is a central claim of the paper, I ask that the code repository be made
available for reviewer inspection.}

\reply The repository is public as of this resubmission --- not on request, and not deferred to
acceptance: \texttt{github.com/helghareeb/certify-feature-reduction} (MIT license). It contains the complete
code, the hash-pinned configurations (the originally submitted configuration and aggregates are
preserved frozen under \texttt{submitted-v1/}), a script that downloads and SHA-256-verifies every
raw dataset, the per-configuration summaries and per-patient caches, the pre-registration with a
dated addendum disclosing every post-registration addition, and a cold-start reproduction guide.
Every number in the revised manuscript regenerates from committed code, the configuration, and a
recorded seed.

\section*{Reviewer 2}

Thank you for the close read --- we applied essentially all of the suggested wording changes, and
they made the paper clearer. Specifics worth noting: the overloaded $k$ is resolved by spelling out
five-fold cross-validation and reserving $k$ exclusively for the number of retained features (which
the revision made load-bearing); the monotonicity claim is now scoped to the mean trajectory and
names the panels, acknowledging that the per-dataset ECE curves (left panel) are not monotone;
Table~3 is deleted and its content folded into one sentence, as suggested; ``a wash on average'' and
``fairness axes agree'' are rewritten in plain terms; ``compact'' and ``wide'' are now formally
defined at first use ($p\le12$ vs.\ $p\ge16$) and marked in the meta-analysis figure; ``Relation to
prior work'' is retitled ``Relation to established practice''; ``co-move'' $\to$ ``move in
tandem''; ``shuffling'' $\to$ ``permuting''; and the Section-1 and Section-3 line edits are applied
as proposed. One deliberate exception: we kept \emph{abstain} (rather than \emph{discard}) because
it is the standard term in the selective-prediction literature, and instead added a defining clause
at first use (``abstain --- decline to issue a prediction'').

\section*{Beyond the requested changes}

Because the reviews prompted a full rebuild, we took the opportunity to strengthen the paper
further: a strict per-dataset safety-certificate table (Section~4.7); leave-one-dataset-out and
bootstrap robustness for the cross-dataset correlations; the population-scale flip-rate
distribution and its race-group breakdown on the full cohort (decision instability is \emph{not}
concentrated on minority groups in this cohort --- reported as measured); reliability diagrams for
the worked-example cell; a selection-stability analysis whose null result rules out ranking noise
as the mechanism of harm; a boxed six-step practitioner checklist; and the enforcement, in code,
of the hash-gate the manuscript describes (analysis now refuses to combine results from mismatched
configurations). A dated addendum to the deposited pre-registration discloses every
post-registration addition explicitly.

\section*{A second round of strengthening, and three claims we withdrew}

After preparing the response above we continued the audit rather than stopping at the reviewers'
requests, and the additional work changed the paper in ways we want to state plainly --- including
where it went against us.

\subsection*{The concession Reviewer~1's line of questioning implied, now measured}

The submitted Limitations conceded that AURC is computed from a confidence, $|p-0.5|$, that is
itself a function of the predicted probability, and that establishing selective reliability as
\emph{independent} evidence ``would require a confidence estimated outside the predicted probability
altogether --- which we did not run.'' We have now run it. For every out-of-fold point we computed
the mean Euclidean distance to its ten nearest neighbours in the full feature space, standardised on
training-fold statistics only, and used that as the confidence; the model's predicted probability
never enters it.

The result corrects our own headline, and we report it in the Abstract and the Discussion rather
than in a limitation. The mean AURC reduction penalty falls from $+0.0334$ under the
probability-scale confidence to $+0.0116$ under the probability-independent one: \textbf{roughly
two-thirds of the measured selective-reliability harm was a property of the probability scale.} The
remaining third survives, decisively ($p=7\times10^{-103}$ over 2{,}160 cells) and with no
dependence on cohort size, which is what makes it evidence of a genuine loss of error-ordering
rather than a restatement of the calibration result. We also tested the obvious model-based
alternative first and report that it fails structurally: per-tree disagreement correlates with
$|p-0.5|$ at $\rho=-0.991$, because a forest disagrees most where its mean prediction sits near the
decision boundary.

\subsection*{``You did not try a flexible enough recalibrator''}

The submitted version rested the claim that recalibration cannot recover selective reliability on
isotonic regression alone. We added temperature scaling and beta calibration under the same
leakage-safe inner-cross-validation protocol and ordered all five recalibrators by flexibility on one
matched grid. Calibration-error recovery rises \emph{strictly monotonically} with flexibility, from
$+0.0054$ at one parameter to $+0.0193$ for the nonparametric recalibrator (Spearman $+1.00$),
reaching $27.8\,\%$ of the uncalibrated baseline; selective-reliability recovery does not follow, and
its correlation with flexibility runs the opposite way ($-0.90$), with the largest movement worth
$3.1\,\%$ of baseline. The failure to recover is therefore not a limitation of the recalibrator.
Net-benefit recovery does rise monotonically, and we report that rather than describe it as flat ---
its magnitude, $0.9\,\%$ of baseline, is what makes it negligible, not its direction.

\subsection*{A causal test of the mechanism, whose pre-registered prediction failed}

Both reviewers' concerns turn on \emph{when} reduction is safe, so we tested our proposed explanation
causally instead of correlationally, and pre-registered the prediction before computing. If diffuse
signal is what makes reduction harmful, then padding a dataset with pure-noise features --- holding
rows, true features, folds and the \emph{absolute} budget fixed --- should increase the harm.

It did not. Dilution lowered the penalty (cleveland, $+0.054\to+0.006$ at ten-fold padding),
reversed its sign (SPECTF, $+0.004\to-0.037$), or left it unmoved (mammographic,
$+0.086\to+0.081$). \textbf{Reduction defends against noise padding}, because the ranker keeps
recovering the true features while the full model degrades. We report the refutation as measured. It
also sharpens the account: what distinguishes a harmed dataset is not how much noise it contains but
how many features carry its signal, an effective dimensionality that no single concentration index
captures across widths differing by two orders of magnitude.

\subsection*{Three claims withdrawn, and why we are telling you}

A pre-submission audit of the manuscript against its own released results found three statements we
could not defend, and we removed them rather than qualify them.

\begin{enumerate}[leftmargin=1.4em,itemsep=2pt]
\item An exploratory claim that the four high-dimensional datasets are \emph{ordered} from least to
      most harmed by signal concentration. Tested across three rankers and three concentration
      statistics it held in two of nine cases, and in \emph{none} of twenty-one once we corrected the
      statistic for the width incomparability that caused the failure. A concentration column in our
      own released results refuted it. Only the prostate-versus-arcene contrast, which is robust
      across every ranker, budget and statistic, is retained.
\item A sentence in the Discussion recommending that budget decisions be ``guided by the
      feature-count / intrinsic-dimension screen above.'' That screen is precisely what this revision
      withdrew in response to Reviewer~1; the sentence was stale text pointing readers at a criterion
      the same manuscript retracts two sections earlier.
\item The description of net-benefit recovery as flat, corrected above.
\end{enumerate}

We mention these because the alternative --- leaving them in, since no reviewer had asked --- is the
failure mode this protocol exists to prevent. Every number in the manuscript is regenerated from a
committed configuration and a recorded seed, every result file's calibration hash resolves to a
configuration shipped in the repository, and every column in every released results file is either
cited in the paper or declared out of scope with a stated reason.


We believe the revised manuscript is substantially stronger than the original, in large part
because of these reviews, and we hope it now meets the bar for acceptance.

\end{document}
